% GENERATED FILE. Edit canonical lessons, then run npm run generate:books.
% canonical-source-hash: fnv1a64:e0622e099f7fe7a5
% canonical-lessons: TA-C426-cantati, TA-C426-iraiccal, TA-C426-tiral, TA-C426-nirampu, TA-C426-tinaru

\chapter{Crowds and Noise}
\label{ch:ta-crowds-and-noise}
\hlchaptermodality{426}

\begin{chapteropening}
\textbf{By the end of this chapter:} \emph{I can talk about bustle and noise, crowds gathering, places filling up and gasping for air.}

Complete the last lesson of chapter 426: I can talk about bustle and noise, crowds gathering, places filling up and gasping for air.
\end{chapteropening}

\section[\texorpdfstring{\ta{சந்தடி} (cantaṭi)}{சந்தடி (cantaṭi)}]{\ta{சந்தடி} (cantaṭi) — bustle; commotion}

\label{lesson:TA-C426-cantati}



\begin{quote}
Before the new one: say the Tamil for a view, then the Tamil for a canteen.
\end{quote}

\subsection*{You'll want to know: \ta{சந்தடி}}
\textbf{\ta{சந்தடி}} — \emph{cantaṭi} — \textquotedblleft{}bustle; commotion\textquotedblright{}.

\begin{quote}
\ta{காலையில்} \ta{சந்தையில்} \ta{ஒரே} \ta{சந்தடி}.
\end{quote}

\begin{quote}
\emph{kālaiyil cantaiyil orē cantaṭi.}
\end{quote}

\begin{quote}
In the morning there is nothing but bustle in the market.
\end{quote}

\ta{சந்தடி} \ta{இல்லாமல்} means quietly, without fuss.

\begin{cousinweb}
\textbf{\ta{சந்தடி}} (\emph{cantaṭi}), \textquotedblleft{}bustle, noise of a crowd\textquotedblright{}, is a Tamil word whose origin is not settled.
\end{cousinweb}

\begin{grammarlens}[title={What you've built}]
One more word for getting around.
\end{grammarlens}

\subsection*{Your turn}
\begin{itemize}
\raggedright
  \item \emph{Say it:} \emph{cantaṭi}
  \item \emph{Say it:} \emph{cantaṭi}, once more
  \item \emph{Say it:} \emph{kālaiyil cantaiyil orē cantaṭi.} — In the morning there is nothing but bustle in the market.
  \item \emph{Recall:} say the Tamil for a view, then the Tamil for a canteen, then say \emph{cantaṭi} again
\end{itemize}

\begin{tcolorbox}[breakable,colback=teal!4,colframe=teal!35!black,title={Before you move on}]
What does \ta{சந்தடி} mean? (\textquotedblleft{}Bustle; commotion\textquotedblright{}.) Say it once more.
\end{tcolorbox}
\section[\texorpdfstring{\ta{இரைச்சல்} (iraiccal)}{இரைச்சல் (iraiccal)}]{\ta{இரைச்சல்} (iraiccal) — noise; din}

\label{lesson:TA-C426-iraiccal}



\begin{quote}
Before the new one: say the Tamil for a canteen, then the Tamil for bustle.
\end{quote}

\subsection*{You'll want to know: \ta{இரைச்சல்}}
\textbf{\ta{இரைச்சல்}} — \emph{iraiccal} — \textquotedblleft{}noise; din\textquotedblright{}.

\begin{quote}
\ta{போக்குவரத்தின்} \ta{இரைச்சலால்} \ta{தூங்க} \ta{முடியவில்லை}.
\end{quote}

\begin{quote}
\emph{pōkkuvarattiṉ iraiccalāl tūṅka muṭiyavillai.}
\end{quote}

\begin{quote}
I could not sleep because of the noise of the traffic.
\end{quote}

The cause takes -\ta{ஆல்}: \ta{இரைச்சலால்}, because of the noise.

\begin{cousinweb}
\textbf{\ta{இரைச்சல்}} (\emph{iraiccal}) is native Dravidian (DEDR 489), \textquotedblleft{}sound, noise, clamour\textquotedblright{}.
\end{cousinweb}

\begin{grammarlens}[title={What you've built}]
One more word for getting around.
\end{grammarlens}

\subsection*{Your turn}
\begin{itemize}
\raggedright
  \item \emph{Say it:} \emph{iraiccal}
  \item \emph{Say it:} \emph{iraiccal}, once more
  \item \emph{Say it:} \emph{pōkkuvarattiṉ iraiccalāl tūṅka muṭiyavillai.} — I could not sleep because of the noise of the traffic.
  \item \emph{Recall:} say the Tamil for a canteen, then the Tamil for bustle, then say \emph{iraiccal} again
\end{itemize}

\begin{tcolorbox}[breakable,colback=teal!4,colframe=teal!35!black,title={Before you move on}]
What does \ta{இரைச்சல்} mean? (\textquotedblleft{}Noise; din\textquotedblright{}.) Say it once more.
\end{tcolorbox}
\section[\texorpdfstring{\ta{திரள்} (tiraḷ)}{திரள் (tiraḷ)}]{\ta{திரள்} (tiraḷ) — to gather; to throng}

\label{lesson:TA-C426-tiral}



\begin{quote}
Before the new one: say the Tamil for bustle, then the Tamil for noise.
\end{quote}

\subsection*{You'll want to know: \ta{திரள்}}
\textbf{\ta{திரள்}} — \emph{tiraḷ} — \textquotedblleft{}to gather; to throng\textquotedblright{}.

\begin{quote}
\ta{சதுக்கத்தில்} \ta{நிறைய} \ta{மக்கள்} \ta{திரண்டார்கள்}.
\end{quote}

\begin{quote}
\emph{catukkattil niṟaiya makkaḷ tiraṇṭārkaḷ.}
\end{quote}

\begin{quote}
Many people thronged the square.
\end{quote}

Past (it …): \ta{திரண்டது}, \emph{tiraṇṭatu}.

\ta{ள்} meets the past -\ta{ந்த}- as \ta{ண்ட}: \ta{திரண்டார்கள்}, they gathered.

\begin{cousinweb}
\textbf{\ta{திரள்}} (\emph{tiraḷ}) is native Dravidian (DEDR 3245), \textquotedblleft{}to become round, to assemble, to crowd together\textquotedblright{}.
\end{cousinweb}

\begin{grammarlens}[title={What you've built}]
One more word for getting around.
\end{grammarlens}

\subsection*{Your turn}
\begin{itemize}
\raggedright
  \item \emph{Say it:} \emph{tiraḷ}
  \item \emph{Say it:} \emph{tiraḷ}, once more
  \item \emph{Say it:} \emph{catukkattil niṟaiya makkaḷ tiraṇṭārkaḷ.} — Many people thronged the square.
  \item \emph{Recall:} say the Tamil for bustle, then the Tamil for noise, then say \emph{tiraḷ} again
\end{itemize}

\begin{tcolorbox}[breakable,colback=teal!4,colframe=teal!35!black,title={Before you move on}]
What does \ta{திரள்} mean? (\textquotedblleft{}To gather; to throng\textquotedblright{}.) Say it once more.
\end{tcolorbox}
\section[\texorpdfstring{\ta{நிரம்பு} (nirampu)}{நிரம்பு (nirampu)}]{\ta{நிரம்பு} (nirampu) — to fill up; to become full}

\label{lesson:TA-C426-nirampu}



\begin{quote}
Before the new one: say the Tamil for noise, then the Tamil for to gather.
\end{quote}

\subsection*{You'll want to know: \ta{நிரம்பு}}
\textbf{\ta{நிரம்பு}} — \emph{nirampu} — \textquotedblleft{}to fill up; to become full\textquotedblright{}.

\begin{quote}
\ta{சில} \ta{நிமிடங்களில்} \ta{பேருந்து} \ta{நிரம்பியது}.
\end{quote}

\begin{quote}
\emph{cila nimiṭaṅkaḷil pēruntu nirampiyatu.}
\end{quote}

\begin{quote}
In a few minutes the bus filled up.
\end{quote}

Past (it …): \ta{நிரம்பியது}, \emph{nirampiyatu}.

\ta{நிரம்பு} is to become full; \ta{நிரப்பு} is to fill something.

\begin{cousinweb}
\textbf{\ta{நிரம்பு}} (\emph{nirampu}) is native Dravidian (DEDR 3672), \textquotedblleft{}to become full, to be complete\textquotedblright{}, in the same article as \ta{நிரப்பு}, \textquotedblleft{}to fill\textquotedblright{}.
\end{cousinweb}

\begin{grammarlens}[title={What you've built}]
One more word for getting around.
\end{grammarlens}

\subsection*{Your turn}
\begin{itemize}
\raggedright
  \item \emph{Say it:} \emph{nirampu}
  \item \emph{Say it:} \emph{nirampu}, once more
  \item \emph{Say it:} \emph{cila nimiṭaṅkaḷil pēruntu nirampiyatu.} — In a few minutes the bus filled up.
  \item \emph{Recall:} say the Tamil for noise, then the Tamil for to gather, then say \emph{nirampu} again
\end{itemize}

\begin{tcolorbox}[breakable,colback=teal!4,colframe=teal!35!black,title={Before you move on}]
What does \ta{நிரம்பு} mean? (\textquotedblleft{}To fill up; to become full\textquotedblright{}.) Say it once more.
\end{tcolorbox}
\section[\texorpdfstring{\ta{திணறு} (tiṇaṟu)}{திணறு (tiṇaṟu)}]{\ta{திணறு} (tiṇaṟu) — to gasp; to be overwhelmed}

\label{lesson:TA-C426-tinaru}



\begin{quote}
Before the new one: say the Tamil for to gather, then the Tamil for to fill up.
\end{quote}

\subsection*{You'll want to know: \ta{திணறு}}
\textbf{\ta{திணறு}} — \emph{tiṇaṟu} — \textquotedblleft{}to gasp; to be overwhelmed\textquotedblright{}.

\begin{quote}
\ta{நெரிசலான} \ta{பேருந்தில்} \ta{நான்} \ta{திணறினேன்}.
\end{quote}

\begin{quote}
\emph{nericalāṉa pēruntil nāṉ tiṇaṟiṉēṉ.}
\end{quote}

\begin{quote}
In the crowded bus I gasped for air.
\end{quote}

Past (I …): \ta{திணறினேன்}, \emph{tiṇaṟiṉēṉ}.

Also of work or questions: \ta{கேள்விகளுக்குப்} \ta{பதில்} \ta{சொல்லத்} \ta{திணறினான்}.

\begin{cousinweb}
\textbf{\ta{திணறு}} (\emph{tiṇaṟu}), \textquotedblleft{}to gasp, to be overwhelmed\textquotedblright{}, is a Tamil verb whose origin is not settled.
\end{cousinweb}

\begin{grammarlens}[title={What you've built}]
One more word for getting around.
\end{grammarlens}

\subsection*{Your turn}
\begin{itemize}
\raggedright
  \item \emph{Say it:} \emph{tiṇaṟu}
  \item \emph{Say it:} \emph{tiṇaṟu}, once more
  \item \emph{Say it:} \emph{nericalāṉa pēruntil nāṉ tiṇaṟiṉēṉ.} — In the crowded bus I gasped for air.
  \item \emph{Recall:} say the Tamil for to gather, then the Tamil for to fill up, then say \emph{tiṇaṟu} again
\end{itemize}

\begin{tcolorbox}[breakable,colback=teal!4,colframe=teal!35!black,title={Before you move on}]
What does \ta{திணறு} mean? (\textquotedblleft{}To gasp; to be overwhelmed\textquotedblright{}.) Say it once more.
\end{tcolorbox}
